\section{Evaluation}\label{planta:sec:evaluation}

This work evaluates \acronym on two use cases from the literature: (i) the SUAVE exemplar~\cite{suave}, in which an AUV performs pipeline inspection, and (ii) an AGV scenario from Cámara et al.~\cite{CAMARASoftwareArchitectureTask2020}, in which the AGV navigates through an environment composed of corridors and rooms. The evaluation is organized around the following research questions:
\begin{description}
    \item[RQ1 (Feasibility).] Can \acronym be integrated into a ROS~2-based robotic system to perform task and architecture co-adaptation at runtime, and how does its mission performance compare with that of existing managing subsystems?
    \item[RQ2 (Performance efficiency).] How do \acronym's planning time and peak memory usage compare with those of the approach by Cámara et al.~\cite{CAMARASoftwareArchitectureTask2020}, and does \acronym produce plans of comparable quality?
    \item[RQ3 (Scalability).] How does \acronym's planning time evolve as the mission and adaptation spaces grow, and what are the trade-offs between optimal and satisficing planning?
\end{description}
RQ1 is addressed by the implementation of \acronym and its integration with \plansys~(\Cref{planta:sec:implementation}), and by applying \acronym to the SUAVE exemplar and its TACA extension~(\Cref{planta:sec:suave_exp}), which primarily serve as a feasibility evaluation. RQ2 is addressed with the AGV scenario~(\Cref{planta:sec:nav_exp}). RQ3 is addressed by varying the AGV scenario along four dimensions, independently and in combination~(\Cref{planta:sec:scale_exp}). For all scenarios, the size of the planning instance is reported in terms of the number of task and reconfiguration actions, derived predicates, objects, and initial facts, to characterize the complexity of \acronym's planning formulations~(\Cref{planta:tab:pddl_instances}). The complete source code and experimental artifacts are publicly available online for the SUAVE experiments~\cite{gustavo_rezende_2026_23619542} and the navigation experiments~\cite{gustavo_rezende_2026_23647225}.

\begin{table}[tbp]
\centering
\caption{Complexity of the planning instances.
    \emph{Task actions}: PDDL actions encoding each task in the mission.
    \emph{Reconfig.\ actions}: PDDL actions encoding architecture reconfiguration.
    \emph{Objects}: domain constants and problem-instance objects.
    \emph{Init.\ facts}: ground atoms in the initial state.
    \emph{Derived predicates}: predicates generated by \owltopddl from the ontology.
    The scalability instances are reported at the maximum tested parameter value
    (\Sone: 900 nodes, \Stwo: 15 architectural configurations, \Sthree: 4 corridor types,
    \Sfour: 5 mission-level actions, and \Sfive: 90 nodes).}%
\label{planta:tab:pddl_instances}
\begin{tabular}{lccccc}
\toprule
\textbf{Scenario} & \makecell{\textbf{Task}\\\textbf{actions}} & \makecell{\textbf{Reconfig.}\\\textbf{actions}} & \textbf{Objects} & \makecell{\textbf{Init.}\\\textbf{facts}} & \makecell{\textbf{Derived}\\\textbf{predicates}} \\ \midrule
SUAVE          & 3 & 2 & 44  & 115  & 54 \\
SUAVE extended & 4 & 2 & 49  & 125  & 55 \\ \midrule
Navigation     & 1 & 2 & 537 & 3184 & 53 \\
\Sone          & 1 & 2 & 947 & 8878 & 53 \\
\Stwo          & 1 & 2 & 150 & 1012 & 53 \\
\Sthree        & 1 & 2 & 144 & 1561 & 53 \\
\Sfour         & 5 & 2 & 153 & 988  & 53 \\
\Sfive         & 5 & 2 & 173 & 1635 & 53 \\
\bottomrule
\end{tabular}
\end{table}

\subsection{SUAVE exemplar}\label{planta:sec:suave_exp}
\paragraph{Experimental design}
To assess \acronym's feasibility, mission performance, and reaction time, it is applied to the SUAVE exemplar~(see \Cref{planta:sec:running_example}) and its TACA extension.

To support the benchmarking of different adaptation approaches, the SUAVE system is structured into a managing and a managed subsystem. The managing subsystem, responsible for the adaptation logic, can be replaced by any self-adaptive system that complies with the provided interfaces. SUAVE is packaged with three distinct managing subsystems: one that arbitrarily selects configurations, a behavior tree-based one, and  Metacontrol~\cite{bozhinoski2022mros}. Additionally, Silva et al.~\cite{rosa} made a ROSA implementation of SUAVE and its extended version available.

To apply \acronym to SUAVE, this work proposes a \pddltomasys model of the SUAVE exemplar based on the SUAVE \tomasys ontology. The main differences are the addition of new actions, utility rules, and a rule to handle uncertainty \Utwo~(water visibility). First, the actions \emph{search pipeline} and \emph{inspect pipeline}, as well as their functional requirements, are added to the SUAVE ontology. For the extended use case the \emph{recharge battery} action is also included. For example, $\requiresF$ relates \emph{search pipeline} to the functions \emph{maintain motion} and \emph{generate search path}. Then, rules are added to infer the $\FunctionsTomasys$ utility, for example, a rule that instantiates the generic (utility) rule of \Cref{planta:sec:tomasys_pddl} for the \emph{generate search path} function. Finally, the following rule is defined to infer whether the \emph{generate search path} function designs are realizable, given the measured water visibility:
\begin{align*}
    &\FunctionDesign(?fd1) \wedge
    \hasQAestimation(?fd1, ?eqa) \wedge
    \isQAtype(?eqa, water\_visibility) \wedge \\
    &\qaHasValue(?eqa, ?eqav) \wedge
    \QAvalue(?mqa) \wedge
    sameAs(?mqa, obs\_water\_visibility) \wedge \\
    & \qaHasValue(?mqa, ?mqav) \wedge
    \isQAtype(?mqa, water\_visibility) \wedge
    lessThan(?mqav, ?eqav) \\
    &\implies \fdRealisability(?fd1, false)
\end{align*}

To exemplify the output of the model transformation process, \Cref{planta:lst:search_pipeline_managed} shows the final PDDL formulation for the search pipeline action from the SUAVE use case. The SUAVE task actions in \Cref{planta:tab:pddl_instances} include the non-managed action of starting the robot. The extended instance adds one extra PDDL action schema, recharge battery, and one derived predicate that encodes the battery level condition for uncertainty~\Uthree.

\paragraph{Experimental Setup}
To compare \acronym with existing adaptation approaches, it was evaluated on both the original and extended SUAVE use cases. To reduce the risk of bias from a single parameter set, as in~\cite{suave, rosa}, the evaluation used randomized values for the robot's initial position, the water visibility shift (i.e., its initial value), and the timing of thruster failures. Specifically, three experiments were conducted for each use case, each with a distinct parameter configuration as detailed in \Cref{planta:tab:suave_params}.

\begin{table}[tbp]
\centering
\caption{Parameter configuration for the SUAVE experiments, where $U(a,b)$ is a random uniform distribution from $a$ to $b$. The initial position is relative to the Gazebo reference frame, where $x$ is fixed at $-17.0$ and $z$ at $-18.5$. The time limit for each execution is set to $250$ seconds.}%
\label{planta:tab:suave_params}
\resizebox{.95\textwidth}{!}{%
\begin{tabular}{llllllll}
\toprule
\multirow{2}{*}{\textbf{Exp.}} &
  \textbf{Initial position} &
  \multicolumn{4}{l}{\textbf{Water visibility}} &
  \multicolumn{2}{l}{\textbf{Thruster failures}} \\
 &
  \textbf{Y} &
  \textbf{Min (m)} &
  \textbf{Max (m)} &
  \textbf{Period (s)} &
  \textbf{Shift (s)} &
  \textbf{Thruster} &
  \textbf{Time (s)} \\ \midrule
1 & $2.5 + U(-0.5, 0.5)$ & 2.5  & 2.5  & N/A & N/A      & 1 and 3 & $100 + U(-100, 100)$ \\
2 & $2.5 + U(-0.5, 0.5)$ & 1.25 & 3.75 & 120 & $U(0,120)$ & N/A     & N/A                \\
3 & $2.5 + U(-0.5, 0.5)$ & 1.25 & 3.75 & 120 & $U(0,120)$ & 1 and 3 & $100 + U(-100, 100)$ \\ \bottomrule
\end{tabular}%
}
\end{table}

The original use case was executed 30 times with no managing subsystem, as well as with each of the following managing subsystems: a random subsystem, a behavior-tree-based subsystem (BT), Metacontrol, ROSA, and \acronym. The extended scenario, which includes TACA, was also executed 30 times for the BT-based subsystem, ROSA, and \acronym. In both the original and the extended scenarios, Fast Downward was configured to use A* search with the blind heuristic for \acronym. To ensure that outliers were averaged out, the parameters were kept the same for each group of five runs (i.e., runs 1–5, 6–10, etc.) within the 30 runs. For each run, standard SUAVE metrics~\cite{suave} were collected: \emph{pipeline search time}, the time from mission start until the AUV locates the pipeline, and \emph{pipeline distance inspected}, the total length of pipeline inspected during the mission. A third metric, \emph{reaction time}~\cite{rosa}, was also collected: the time elapsed from a constraint violation until the managing system completes adaptation of the managed subsystem.

\paragraph{Results}
The experimental results for the SUAVE use case are presented in \Cref{planta:tab:suave_results} and for the extended use case in \Cref{planta:tab:suave_extended_results}. To verify whether there is a statistically significant difference between the results, the Wilcoxon signed-rank test was applied to all pairs of managing subsystems for both the \emph{pipeline search time} and \emph{pipeline distance inspected} metrics. The p-values for experiment 3 and extended experiment 3, which combine thruster failures and water visibility changes, are reported in \Cref{planta:tab:wilcoxon}. The p-values for the remaining experiments are reported in the electronic supplement.

Results were paired by run index, and pairs in which either system did not find the pipeline were excluded from both comparisons. For each metric, the alternative hypothesis is that the managing system in the row achieves a lower \emph{pipeline search time} or, respectively, a greater \emph{pipeline distance inspected} than the one in the column.
p-values less than 0.05 indicate statistical significance in favor of the alternative hypothesis.

% Generated by scripts/summarize_suave_results.py
% Source campaigns: /tmp/claude-1000/-home-gus-ros-workspaces-planta-ws/90f03437-a284-4a9e-9032-238a88aa72ef/scratchpad/paper_tables_input/campaigns
% Do not edit by hand.
\begin{table}[tb]
\centering
\caption{SUAVE results. The Random managing subsystem found the pipeline in 21/30 runs (70.00\%) for Experiment 1, 22/30 runs (73.33\%) for Experiment 2, and 22/30 runs (73.33\%) for Experiment 3. All other managing subsystems found the pipeline in every run. Search time and inspected distance report the mean and sample standard deviation only for runs that found the pipeline. \Uone denotes thruster failure and \Utwo water visibility. Reaction times are arithmetic means of recorded events per uncertainty across all retained runs. N/A denotes an inapplicable metric or insufficient observations. The initial visibility adjustment in Experiment 1 is excluded.}\label{planta:tab:suave_results}
\resizebox{.95\textwidth}{!}{%
\begin{tabular}{llcccccc}
\toprule
\multirow{2}{*}{\textbf{Exp.}} & \multirow{2}{*}{\textbf{Managing subsystem}} & \multicolumn{2}{c}{\textbf{Search time (s)}} & \multicolumn{2}{c}{\textbf{Distance inspected (m)}} & \multicolumn{2}{c}{\textbf{Mean reaction time (s)}} \\
 & & \textbf{Mean} & \textbf{Std.\ dev.} & \textbf{Mean} & \textbf{Std.\ dev.} & \textbf{\Uone} & \textbf{\Utwo} \\
\midrule
\multirow{6}{*}{1} & None & 58.77 & 26.07 & 49.50 & 12.56 & N/A & N/A \\
 & Random & 126.71 & 52.90 & 11.65 & 10.08 & N/A & N/A \\
 & BT & 59.63 & 27.46 & 58.17 & 8.43 & 0.07 & N/A \\
 & Metacontrol & 59.20 & 25.50 & 58.92 & 8.63 & 1.11 & N/A \\
 & ROSA & 65.33 & 27.48 & 57.76 & 9.06 & 0.75 & N/A \\
 & \textbf{\acronym} & \textbf{65.07} & \textbf{33.95} & \textbf{43.56} & \textbf{8.77} & \textbf{0.88} & \textbf{N/A} \\
\midrule
\multirow{6}{*}{2} & None & 152.13 & 30.32 & 27.31 & 11.23 & N/A & N/A \\
 & Random & 113.32 & 48.30 & 10.91 & 8.15 & N/A & N/A \\
 & BT & 83.83 & 30.96 & 50.42 & 10.25 & N/A & 0.13 \\
 & Metacontrol & 85.53 & 31.51 & 50.22 & 10.37 & N/A & 1.37 \\
 & ROSA & 89.03 & 31.19 & 49.69 & 10.46 & N/A & 1.52 \\
 & \textbf{\acronym} & \textbf{87.83} & \textbf{31.41} & \textbf{44.45} & \textbf{10.11} & \textbf{N/A} & \textbf{1.78} \\
\midrule
\multirow{6}{*}{3} & None & 148.97 & 33.14 & 25.90 & 9.47 & N/A & N/A \\
 & Random & 119.50 & 52.35 & 8.76 & 7.31 & N/A & N/A \\
 & BT & 84.43 & 30.94 & 49.82 & 10.30 & 0.07 & 0.12 \\
 & Metacontrol & 86.77 & 31.48 & 49.30 & 10.08 & 1.22 & 1.52 \\
 & ROSA & 89.57 & 31.71 & 49.31 & 10.53 & 0.78 & 1.56 \\
 & \textbf{\acronym} & \textbf{88.73} & \textbf{32.07} & \textbf{43.58} & \textbf{9.45} & \textbf{1.45} & \textbf{1.86} \\
\bottomrule
\end{tabular}%
}
\end{table}


% Generated by scripts/summarize_suave_results.py
% Source campaigns: /tmp/claude-1000/-home-gus-ros-workspaces-planta-ws/90f03437-a284-4a9e-9032-238a88aa72ef/scratchpad/paper_tables_input/campaigns
% Do not edit by hand.
\begin{table}[tb]
\centering
\caption{SUAVE extended results. All managing subsystems found the pipeline in every run. Search time and inspected distance report the mean and sample standard deviation only for runs that found the pipeline. \Uone denotes thruster failure and \Utwo water visibility, with \Uthree denoting battery level. Reaction times are arithmetic means of recorded events per uncertainty across all retained runs. N/A denotes an inapplicable metric or insufficient observations. The initial visibility adjustment in Experiment 1 is excluded.}\label{planta:tab:suave_extended_results}
\resizebox{.95\textwidth}{!}{%
\begin{tabular}{llccccccc}
\toprule
\multirow{2}{*}{\textbf{Exp.}} & \multirow{2}{*}{\textbf{Managing subsystem}} & \multicolumn{2}{c}{\textbf{Search time (s)}} & \multicolumn{2}{c}{\textbf{Distance inspected (m)}} & \multicolumn{3}{c}{\textbf{Mean reaction time (s)}} \\
 & & \textbf{Mean} & \textbf{Std.\ dev.} & \textbf{Mean} & \textbf{Std.\ dev.} & \textbf{\Uone} & \textbf{\Utwo} & \textbf{\Uthree} \\
\midrule
\multirow{3}{*}{1} & BT & 64.30 & 34.37 & 23.40 & 8.17 & 0.09 & N/A & 1.09 \\
 & ROSA & 61.10 & 24.85 & 22.41 & 6.95 & 0.81 & N/A & 1.95 \\
 & \textbf{\acronym} & \textbf{71.47} & \textbf{37.23} & \textbf{17.47} & \textbf{8.79} & \textbf{0.89} & \textbf{N/A} & \textbf{1.44} \\
\midrule
\multirow{3}{*}{2} & BT & 83.93 & 30.88 & 16.74 & 9.09 & N/A & 0.12 & 1.08 \\
 & ROSA & 88.70 & 31.50 & 15.40 & 8.96 & N/A & 1.57 & 2.03 \\
 & \textbf{\acronym} & \textbf{88.20} & \textbf{31.86} & \textbf{14.72} & \textbf{8.44} & \textbf{N/A} & \textbf{2.37} & \textbf{2.25} \\
\midrule
\multirow{3}{*}{3} & BT & 89.07 & 40.14 & 16.25 & 9.07 & 0.07 & 0.12 & 1.07 \\
 & ROSA & 88.10 & 31.71 & 15.48 & 8.96 & 0.76 & 1.67 & 1.96 \\
 & \textbf{\acronym} & \textbf{89.30} & \textbf{32.06} & \textbf{14.29} & \textbf{8.47} & \textbf{1.57} & \textbf{2.40} & \textbf{2.31} \\
\bottomrule
\end{tabular}%
}
\end{table}


\begin{table}[tbp]
\centering
\caption{Paired Wilcoxon signed-rank comparisons for experiment 3 and extended experiment 3. Each cell reports the unadjusted p-values for the comparison between the managing subsystem in the row and the one in the column for both the \emph{pipeline search time} and \emph{pipeline distance inspected} metrics. Cells with a green background indicate that the alternative hypothesis was favored, and cells with a blue background indicate that the null hypothesis cannot be rejected in favor of the alternative hypothesis. Values rounding to 0.000 are shown as $<0.001$. ``MC'' stands for Metacontrol.}%
\label{planta:tab:wilcoxon}
\begin{subtable}{\textwidth}
\centering
\inputtablebody{pddl/tables/wilcoxon_analysis/exp3_wilcoxon.tex}
\caption{Experiment 3}\label{planta:tab:wilcoxon-exp3}
\end{subtable}

\medskip
% 0.78 matches the font size of (a): the generated table scales to 0.7\textwidth
\begin{subtable}{0.78\textwidth}
\centering
\inputtablebody{pddl/tables/wilcoxon_analysis/extended_exp3_wilcoxon.tex}
\caption{Extended experiment 3}\label{planta:tab:wilcoxon-extended-exp3}
\end{subtable}
\end{table}

\emph{SUAVE experiments:} Across the three SUAVE experiments (\Cref{planta:tab:wilcoxon-exp3} and \Cref{planta:S-tab:wilcoxon-exp1,planta:S-tab:wilcoxon-exp2}), the BT managing subsystem achieved significantly lower \emph{pipeline search times} than all other configurations. Metacontrol achieved significantly lower search times than ROSA and \acronym in every experiment. \acronym achieved significantly lower search times than ROSA in experiment 2, while no statistically significant difference between them was detected in experiments 1 and 3. For \emph{pipeline distance inspected}, BT, Metacontrol, and ROSA each achieved significantly greater distances than \acronym in all three experiments. \acronym achieved significantly greater inspected distances than Random in every experiment. Compared with the configuration without a managing subsystem, \acronym inspected significantly less pipeline distance in experiment 1 and significantly more in experiments 2 and 3. % chktex 12 -- sentence starts with a macro

\emph{SUAVE extended experiments:} Across the three extended SUAVE experiments (\Cref{planta:tab:wilcoxon-extended-exp3} and \Cref{planta:S-tab:wilcoxon-extended-exp1,planta:S-tab:wilcoxon-extended-exp2}), the BT managing subsystem achieved significantly lower \emph{pipeline search times} than ROSA and \acronym, while no statistically significant difference in search time was detected between ROSA and \acronym. For \emph{pipeline distance inspected}, BT achieved significantly greater distances than both ROSA and \acronym, and ROSA achieved significantly greater distances than \acronym in every experiment.

\emph{Answer to RQ1:} These results support the feasibility of using \acronym for runtime task and architecture co-adaptation in SUAVE, but also reveal limitations in mission performance. BT, Metacontrol, and ROSA each achieved significantly better results than \acronym on at least one metric in every experiment in which they were evaluated. \acronym performed significantly better than Random on both metrics in all original experiments, and better than the configuration without a managing subsystem in experiments 2 and 3, but worse in experiment 1. BT's shorter mean reaction times (\Cref{planta:tab:suave_results,planta:tab:suave_extended_results}) may partly explain its advantage, although the experiments do not establish a causal relationship between reaction time and mission performance. % chktex 12 -- sentence starts with a macro

\subsection{Navigation use case}\label{planta:sec:nav_exp}
\paragraph{Experimental design}
To evaluate \acronym's performance efficiency, in terms of planning time and peak resident set size (RSS) memory usage, and its plan quality, it was applied to the navigation scenario from~\cite{CAMARASoftwareArchitectureTask2020}. Results were compared to those obtained from a replicated implementation of the algorithm proposed in~\cite{CAMARASoftwareArchitectureTask2020}, hereafter referred to as \CamaraApproach.

In the scenario, a robot must navigate from an initial position to a goal position within an environment composed of corridors and rooms. Some corridors might be dark or contain obstacles, reducing their safety. During navigation, the robot dynamically adapts both its architecture (i.e., the sensors and localization algorithms employed) and its task execution (i.e., the path followed) to circumvent failures or prioritize safety or energy consumption. For example, if the robot's planned path includes a dark corridor and its lamp fails during execution, it must reconfigure its architecture and adapt its task execution to avoid dark corridors.

\paragraph{Experimental setup}

To compare \acronym's time performance and memory usage with \CamaraApproach, the map from~\cite{CAMARASoftwareArchitectureTask2020} was discretized to ensure uniform edge lengths. This change was necessary because \acronym does not support numeric fluents and therefore cannot handle edges of varying sizes~(\Cref{planta:sec:limitations}).
\acronym was then applied to perform task and configuration planning on the discretized map. For this comparison, Fast Downward was configured to use eager greedy search with the goal-count heuristic.

Additionally, a model generator was developed to produce PRISM models equivalent to those used in the final step of \CamaraApproach\footnote{The complete algorithm from~\cite{CAMARASoftwareArchitectureTask2020} could not be replicated.}, enabling policy generation for task planning and reconfiguration. A minor correction was introduced to the original PRISM model by adding an effect to the move actions, allowing the selection of distinct configurations for each corridor.  This work assumes that the correction aligns with the authors’ original intent, as the published model appears to contain a typographical error. For all experiments, PRISM was executed with the goal of maximizing the amount of battery remaining at the end of the mission. Both approaches were evaluated by applying them ten times to all possible combinations of initial and goal locations.

\paragraph{Results}
\Cref{planta:fig:camara_pddl_sas} shows that, at the median, \acronym generates both task and configuration plans $3.00$ times faster than \CamaraApproach generates the corresponding policy from the PRISM models, with $6.84$ times smaller peak memory usage. In the scenarios with the worst planning time, \acronym (from location 16 to 25) takes $0.56$ seconds with $66.88$~MiB peak memory usage, compared to $13.71$ seconds with $4089.03$~MiB peak memory usage for \CamaraApproach (from location 52 to 24). Notably, the spread of execution times and memory consumption is much higher for \CamaraApproach, making the resources required by \acronym more predictable. These results should be interpreted as a comparison of the computational cost of two different ways of addressing TACA rather than as a direct comparison of approaches that pursue the same objective. \CamaraApproach seeks an optimal task and configuration policy under probabilistic uncertainty, whereas \acronym favors satisficing planning over mission-level optimality, planning from the currently observed state and handling changes during execution through replanning. % chktex 12 -- sentence starts with a macro

Despite the difference in planning time and memory usage, the two approaches produce plans with almost identical path length and remaining battery. The distributions of distance traveled and of the numbers of move and reconfiguration actions nearly coincide. For example, the interquartile range of distance traveled is $36.81$ to $82.70$ for \acronym and $36.62$ to $82.66$ for \CamaraApproach, with the same median of $59.89$. The median remaining battery differs by only $1.02\%$ ($25219.50$ for \acronym against $25480.00$ for \CamaraApproach). Both approaches reconfigure more than once in a minority of the scenarios, $11.85\%$ for \acronym and $15.52\%$ for \CamaraApproach, and up to four times (for example, between locations 11 and 17 for \acronym and between locations 5 and 18 for \CamaraApproach). % chktex 12 -- sentence starts with a macro

\begin{figure}[tbp]
    \centering
    \hfill
    \begin{subfigure}[t]{0.45\linewidth}
        \vspace{0pt} % Aligns top
        \centering
        % Generated by navigation_planta.paper_export; n=30800; Tukey 1.5 IQR; 2731 outliers, 250 rank-sampled marks.
\def\camaraPrismTimeLW{1.154856}
\def\camaraPrismTimeLQ{1.208374}
\def\camaraPrismTimeMED{1.212604}
\def\camaraPrismTimeUQ{2.516190}
\def\camaraPrismTimeUW{4.449019}
%
% Generated by navigation_planta.paper_export; n=30800; Tukey 1.5 IQR; 3964 outliers, 250 rank-sampled marks.
\def\camaraPddlTimeLW{0.402322}
\def\camaraPddlTimeLQ{0.403349}
\def\camaraPddlTimeMED{0.403796}
\def\camaraPddlTimeUQ{0.404449}
\def\camaraPddlTimeUW{0.406096}
%
% Generated by navigation_planta; do not edit by hand.
\begin{tikzpicture}[baseline=(paperAxis.north)]
\begin{axis}[name=paperAxis,
width=\textwidth,
height=0.75\textwidth,
boxplot/draw direction=y,
xmin=0,
xmax=3,
xtick={1,2},
xticklabels={\CamaraApproach,\paperEagerLabel},
ylabel={Planning time [s]},
ymajorgrids=true,
ymode=log,
ymin=0.3,
ymax=30,
ytick={1,10},
extra y ticks={0.5,30},
extra y tick labels={0.5,30}
]
\addplot+[boxplot prepared={lower whisker=\camaraPrismTimeLW,lower quartile=\camaraPrismTimeLQ,median=\camaraPrismTimeMED,upper quartile=\camaraPrismTimeUQ,upper whisker=\camaraPrismTimeUW,draw position=1,box extend=0.2},draw=\blueColor,fill=\blueColor!30,line width=0.75pt,mark=o,mark options={fill=\blueColor,draw=\blueColor,opacity=0.2,line width=0.5pt}] table[y=y,col sep=comma] {pddl/navigation_experiment/data/camara_map_prism_experiment_time_outliers.csv}
node[left,text=\blueColor] at (boxplot box cs: \camaraPrismTimeMED,0) {\pgfmathprintnumber{\camaraPrismTimeMED}};
\addplot+[boxplot prepared={lower whisker=\camaraPddlTimeLW,lower quartile=\camaraPddlTimeLQ,median=\camaraPddlTimeMED,upper quartile=\camaraPddlTimeUQ,upper whisker=\camaraPddlTimeUW,draw position=2,box extend=0.2},draw=\greenColor,fill=\greenColor!30,line width=0.75pt,mark=o,mark options={fill=\greenColor,draw=\greenColor,opacity=0.2,line width=0.5pt}] table[y=y,col sep=comma] {pddl/navigation_experiment/data/camara_map_pddl_experiment_time_outliers.csv}
node[left,text=\greenColor] at (boxplot box cs: \camaraPddlTimeMED,0) {\pgfmathprintnumber{\camaraPddlTimeMED}};
\end{axis}
\end{tikzpicture}

        \caption{Planning time comparison between \acronym and \CamaraApproach.}%
        \label{planta:fig:time_comparison}
    \end{subfigure}
    \hfill
    \begin{subfigure}[t]{0.45\linewidth}
        \vspace{0pt} % Aligns top
        \centering
        % Generated by navigation_planta.paper_export; n=30800; Tukey 1.5 IQR; 1091 outliers, 250 rank-sampled marks.
\def\camaraPrismMemLW{111.222656}
\def\camaraPrismMemLQ{190.210938}
\def\camaraPrismMemMED{461.085938}
\def\camaraPrismMemUQ{1147.773437}
\def\camaraPrismMemUW{2584.039062}
%
% Generated by navigation_planta.paper_export; n=30800; Tukey 1.5 IQR; 307 outliers, 250 rank-sampled marks.
\def\camaraPddlMemLW{64.460938}
\def\camaraPddlMemLQ{66.527344}
\def\camaraPddlMemMED{67.390625}
\def\camaraPddlMemUQ{67.906250}
\def\camaraPddlMemUW{68.378906}
%
% Generated by navigation_planta; do not edit by hand.
\begin{tikzpicture}[baseline=(paperAxis.north)]
\begin{axis}[name=paperAxis,
width=\textwidth,
height=0.75\textwidth,
boxplot/draw direction=y,
xmin=0,
xmax=3,
xtick={1,2},
xticklabels={\CamaraApproach,\paperEagerLabel},
ylabel={RSS Peak Memory [MiB]},
ymajorgrids=true,
ymode=log,
ymin=5,
ymax=10000,
ytick={10,100,1000,10000}
]
\addplot+[boxplot prepared={lower whisker=\camaraPrismMemLW,lower quartile=\camaraPrismMemLQ,median=\camaraPrismMemMED,upper quartile=\camaraPrismMemUQ,upper whisker=\camaraPrismMemUW,draw position=1,box extend=0.2},draw=\blueColor,fill=\blueColor!30,line width=0.75pt,mark=o,mark options={fill=\blueColor,draw=\blueColor,opacity=0.2,line width=0.5pt}] table[y=y,col sep=comma] {pddl/navigation_experiment/data/camara_map_prism_experiment_peak_memory_outliers.csv}
node[left,text=\blueColor] at (boxplot box cs: \camaraPrismMemMED,0) {\pgfmathprintnumber{\camaraPrismMemMED}};
\addplot+[boxplot prepared={lower whisker=\camaraPddlMemLW,lower quartile=\camaraPddlMemLQ,median=\camaraPddlMemMED,upper quartile=\camaraPddlMemUQ,upper whisker=\camaraPddlMemUW,draw position=2,box extend=0.2},draw=\greenColor,fill=\greenColor!30,line width=0.75pt,mark=o,mark options={fill=\greenColor,draw=\greenColor,opacity=0.2,line width=0.5pt}] table[y=y,col sep=comma] {pddl/navigation_experiment/data/camara_map_pddl_experiment_peak_memory_outliers.csv}
node[left,text=\greenColor] at (boxplot box cs: \camaraPddlMemMED,0) {\pgfmathprintnumber{\camaraPddlMemMED}};
\end{axis}
\end{tikzpicture}

        \caption{Peak RSS memory comparison between \acronym and \CamaraApproach.}%
        \label{planta:fig:memory_comparison}
    \end{subfigure}
    \hfill
    \caption{Comparison between \acronym and \CamaraApproach. The number is the median value, the boxes show interquartile range (Q1--Q3), the whiskers reach until $\pm 1.5$~IQR, and the circles show a sample of the observations beyond the whiskers.}%
    \label{planta:fig:camara_pddl_sas}
\Description{Two box-plot panels compare the replicated approach \CamaraApproach with \acronym. The
left panel shows planning time in seconds, with medians of 1.21 for the replicated
approach and 0.40 for \acronym. The right panel shows peak memory in MiB, with
medians of 461.09 and 67.39, respectively. The replicated approach has a wider spread
and larger high-value outliers in both panels.}
\end{figure}

The reported computational performance difference likely underestimates \acronym's advantage. In~\cite{CAMARASoftwareArchitectureTask2020}, the final step, the sole component that could be reproduced for \CamaraApproach, accounts for only 39.34\% of the total execution time. Furthermore, \CamaraApproach generates and verifies PRISM models for the ten best paths between the initial and goal locations, selecting the best policy according to a predefined utility function. In contrast, the experiment described in this paper generates and verifies a PRISM model only for the single best path.

\emph{Answer to RQ2:} In the evaluated AGV scenario, \acronym addresses TACA at a lower computational cost than the replicated final step of \CamaraApproach, in both planning time and peak memory usage, while producing plans with almost identical path length and remaining battery.

\subsection{Scalability}\label{planta:sec:scale_exp}

\paragraph{Experimental design} To assess the scalability of the proposed approach with respect to mission and adaptation spaces, in terms of planning time and memory usage, five additional experiments were conducted using variations of the previous navigation scenario. Specifically, four experiments were designed to independently evaluate the effect of varying the map size, the number of architectural configurations, the number of quality requirements, and the number of task actions. A fifth experiment was designed to stress-test the approach by varying the map size while fixing all other parameters at their maximum values.

Furthermore, to isolate the effect of the adaptation logic on planning time and memory usage, each of the scalability experiments described above was also conducted using only the mission logic (i.e., without reconfiguration). In addition, to compare optimal and satisficing planning, all scalability experiments, both with and without the adaptation logic, were repeated with Fast Downward configured to use A* search with the blind heuristic and eager greedy search with the goal-count heuristic, where an optimal plan corresponds to the shortest plan since all actions have unit cost.

\paragraph{Experimental setup}

For each scalability experiment, a map generator produces grid-like maps with the required number of nodes, edges, and corridor types. The generator skips $10\%$ of nodes and $15\%$ of the edges are removed while preserving graph connectivity. In the baseline maps, $25\%$ of the corridors are marked as dark and $25\%$ are marked as unsafe. The generated map is then translated into a PDDL problem file, leveraging the Unified Planning~\cite{unified_planning_2025_micheli} library, with node $0$ as the robot's initial position and the highest-indexed node as the goal waypoint. \Cref{planta:fig:grid_map_example} shows an example map. Each experiment is repeated $10$ times for each search strategy~(i.e., A* and eager greedy) and adaptation mode~(with and without adaptation), while the \emph{planning time}, \emph{peak memory usage}, and the \emph{number of actions} in the plan are collected.

\begin{figure}[tbp]
    \centering
    \resizebox{0.30\textwidth}{!}{
        \begin{tikzpicture}[font=\sffamily\bfseries]
    
    % Create a 5x5 grid of nodes
    \foreach \i in {0,1,2,3,4,5,6} {
        \foreach \j in {0,1,2,3,4,5,6} {
            % Name each node with its coordinates
            \node at (\i, -0.75*\j) (node-\i-\j) {};
            \node[gridNode] at (\i, -0.75*\j)  {};
        }
    }
    
    \node[goal] at (0, 0.75) (goal) {G};
    
    \node[start] at (node-0-6) (start) {I};
    
    % hide deleted nodes
    
    \node[hideNode] at (node-2-3) {};
    \node[hideNode] at (node-3-4) {};
    \node[hideNode] at (node-3-6) {};
    \node[hideNode] at (node-6-2) {};
    \node[hideNode] at (node-6-5) {};
    
    
    \begin{pgfonlayer}{background}
        % colums
        \draw[corridor] (goal.center) -- (node-0-0);
        
        \draw[corridor] (node-0-0) -- (node-0-1);
        \draw[corridor] (node-0-1) -- (node-0-2);
        \draw[corridor] (node-0-2) -- (node-0-3);
        %\draw[corridor] (node-0-3) -- (node-0-4);
        \draw[corridorUnsafe] (node-0-4) -- (node-0-5);
        \draw[corridorUnsafe] (node-0-5) -- (node-0-6);
        
        %\draw[corridor] (node-1-0) -- (node-1-1);
        \draw[corridor] (node-1-1) -- (node-1-2);
        \draw[corridor] (node-1-2) -- (node-1-3);
        \draw[corridorDarkUnsafe] (node-1-3) -- (node-1-4);
        \draw[corridorDarkUnsafe] (node-1-4) -- (node-1-5);
        \draw[corridor] (node-1-5) -- (node-1-6);
        
        %\draw[corridor] (node-2-0) -- (node-2-1);
        \draw[corridor] (node-2-1) -- (node-2-2);
        %\draw[corridor] (node-2-2) -- (node-2-3);
        %\draw[corridor] (node-2-3) -- (node-2-4);
        \draw[corridorUnsafe] (node-2-4) -- (node-2-5);
        \draw[corridorDark] (node-2-5) -- (node-2-6);
        
        \draw[corridorUnsafe] (node-3-0) -- (node-3-1);
        %\draw[corridor] (node-3-1) -- (node-3-2);
        %\draw[corridor] (node-3-2) -- (node-3-3);
        %\draw[corridor] (node-3-3) -- (node-3-4);
        %\draw[corridor] (node-3-4) -- (node-3-5);
        %\draw[corridor] (node-3-5) -- (node-3-6);
        
        \draw[corridor] (node-4-0) -- (node-4-1);
        \draw[corridor] (node-4-1) -- (node-4-2);
        \draw[corridor] (node-4-2) -- (node-4-3);
        \draw[corridor] (node-4-3) -- (node-4-4);
        \draw[corridorUnsafe] (node-4-4) -- (node-4-5);
        \draw[corridor] (node-4-5) -- (node-4-6);
        
        \draw[corridorDark] (node-5-0) -- (node-5-1);
        %\draw[corridor] (node-5-1) -- (node-5-2);
        \draw[corridor] (node-5-2) -- (node-5-3);
        \draw[corridorDark] (node-5-3) -- (node-5-4);
        \draw[corridorUnsafe] (node-5-4) -- (node-5-5);
        \draw[corridor] (node-5-5) -- (node-5-6);
        
        \draw[corridorDark] (node-6-0) -- (node-6-1);
        %\draw[corridor] (node-6-1) -- (node-6-2);
        %\draw[corridor] (node-6-2) -- (node-6-3);
        \draw[corridor] (node-6-3) -- (node-6-4);
        %\draw[corridor] (node-6-4) -- (node-6-5);
        %\draw[corridor] (node-6-5) -- (node-6-6);
        
        % rows
        \draw[corridorDark] (node-0-0) -- (node-1-0);
        \draw[corridorUnsafe] (node-1-0) -- (node-2-0);
        \draw[corridor] (node-2-0) -- (node-3-0);
        \draw[corridorDark] (node-3-0) -- (node-4-0);
        \draw[corridor] (node-4-0) -- (node-5-0);
        \draw[corridor] (node-5-0) -- (node-6-0);
        
        \draw[corridor] (node-0-1) -- (node-1-1);
        \draw[corridorUnsafe] (node-1-1) -- (node-2-1);
        %\draw[corridor] (node-2-1) -- (node-3-1);
        \draw[corridor] (node-3-1) -- (node-4-1);
        \draw[corridor] (node-4-1) -- (node-5-1);
        \draw[corridorUnsafe] (node-5-1) -- (node-6-1);
        
        \draw[corridorUnsafe] (node-0-2) -- (node-1-2);
        \draw[corridorDark] (node-1-2) -- (node-2-2);
        \draw[corridor] (node-2-2) -- (node-3-2);
        \draw[corridor] (node-3-2) -- (node-4-2);
        \draw[corridorDark] (node-4-2) -- (node-5-2);
        %\draw[corridor] (node-5-2) -- (node-6-2);
        
        \draw[corridorDark] (node-0-3) -- (node-1-3);
        %\draw[corridor] (node-1-3) -- (node-2-3);
        %\draw[corridor] (node-2-3) -- (node-3-3);
        \draw[corridor] (node-3-3) -- (node-4-3);
        \draw[corridor] (node-4-3) -- (node-5-3);
        %\draw[corridor] (node-5-3) -- (node-6-3);
        
        %\draw[corridor] (node-0-4) -- (node-1-4);
        \draw[corridor] (node-1-4) -- (node-2-4);
        %\draw[corridor] (node-2-4) -- (node-3-4);
        %\draw[corridor] (node-3-4) -- (node-4-4);
        \draw[corridor] (node-4-4) -- (node-5-4);
        \draw[corridorUnsafe] (node-5-4) -- (node-6-4);
        
        \draw[corridor] (node-0-5) -- (node-1-5);
        \draw[corridorDark] (node-1-5) -- (node-2-5);
        \draw[corridor] (node-2-5) -- (node-3-5);
        %\draw[corridor] (node-3-5) -- (node-4-5);
        \draw[corridorUnsafe] (node-4-5) -- (node-5-5);
        %\draw[corridor] (node-5-5) -- (node-6-5);
        
        \draw[corridor] (node-0-6) -- (node-1-6);
        \draw[corridorDark] (node-1-6) -- (node-2-6);
        %\draw[corridor] (node-2-6) -- (node-3-6);
        %\draw[corridor] (node-3-6) -- (node-4-6);
        \draw[corridor] (node-4-6) -- (node-5-6);
        \draw[corridorDark] (node-5-6) -- (node-6-6);
   \end{pgfonlayer}
    
\end{tikzpicture}
    }
    \caption{Grid map example with $45$ retained nodes out of $50$ generated nodes.
    White edges represent normal corridors, filled edges represent dark corridors and curved edges represent unsafe corridors. The initial location is marked with {\sffamily\bfseries I} and the goal location is marked with {\sffamily\bfseries G}.
    }%
    \label{planta:fig:grid_map_example}
\Description{Grid-like navigation graph with 45 locations, consisting of a seven-by-seven grid with five locations removed and an additional goal location. Several connections are omitted. The
initial location I is at the lower left, and the goal G is above the upper-left grid
location. Edges connect neighboring locations: white straight edges denote normal
corridors, filled edges denote dark corridors, and curved edges denote unsafe corridors.
Filled curved edges denote corridors that are both dark and unsafe.}
\end{figure}

\Cref{planta:tab:scalability_params} lists the parameter ranges of the five scalability experiments. Architectural configurations are combinations of sensors and localization algorithms, and \Stwo adds further combinations\footnote{The additional architectural configurations are included to increase the size of the architectural configuration space, not to model realistic selection criteria or quality-attribute trade-offs.}. \Sthree models additional quality requirements as two new corridor types: outdoor corridors ($20\%$) and dusty corridors ($20\%$). Outdoor corridors require sensors that are robust to sunlight, for example, lidars, whereas dusty corridors require sensors that are robust to floating particles, for example, RGB cameras. \Sfour adds the mission-level PDDL actions inspect, deliver, recharge, and report. The map-size upper bound in \Sfive is lower than in \Sone because \Sfive combines the largest architectural configuration space, the largest set of corridor requirements, and the largest mission-action set, which substantially increases the planning space.

\begin{table}[tbp]
  \centering
  \caption{Parameter configurations for the five scalability experiments~(S1--S5).
    Map nodes are generated nodes. The map generator skips $10\%$ of them, so the maximum of $1000$ generated nodes in S1 and $100$ in S5 corresponds to $900$ and $90$ retained nodes, respectively.
    For S4 and S5, the number of architectural configurations grows with the number of action types: each additional mission-level PDDL action introduces two further configurations.
  }%
  \label{planta:tab:scalability_params}
  \begin{tabular}{llcccc}
    \toprule
    & \textbf{Experiment} &
    \makecell{\textbf{Map nodes}} &
    \makecell{\textbf{Arch.}\\\textbf{configs}} &
    \makecell{\textbf{Corridor}\\\textbf{types}} &
    \makecell{\textbf{Actions}} \\
    \midrule
    \Sone & Map size               & $10$--$1000$, step $10$ & $6$                 & $2$     & $1$ \\
    \Stwo & Arch.\ configurations  & $100$                   & $2,4,6,8,10,12,15$  & $2$     & $1$ \\
    \Sthree & Quality requirements   & $100$                   & $6$                 & $2,3,4$ & $1$ \\
    \Sfour & Mission-level actions           & $100$                   & $6$--$14$           & $2$     & $1,2,3,4,5$ \\
    \Sfive & Combined (stress test) & $10$--$100$, step $10$  & $23$                & $4$     & $5$ \\
    \bottomrule
  \end{tabular}
\end{table}

Experiments \Stwo, \Sthree, and \Sfour use maps with $100$ nodes, so that changes in planning time and memory usage can be interpreted primarily with respect to the dimension varied in each experiment rather than map size, which is evaluated separately in \Sone. \Sfive combines the action-driven configuration growth of \Sfour with the additional architectural configurations introduced in \Stwo, yielding $23$ configurations. The configurations added in \Sfour are intentionally not ranked by utility, which avoids introducing a separate utility-optimization factor into the mission-action scalability experiment. \Cref{planta:tab:pddl_instances} reports the resulting complexity of each experiment's planning instance at its maximum tested parameter value.

\paragraph{Results}

Across the five scalability experiments, the number of mission-level PDDL actions is the dominant driver of planning time, with map size as a secondary driver. Architectural configurations (\Stwo) and quality requirements (\Sthree) have a small effect on planning time. Map size (\Sone) produces a moderate increase for both search strategies, whereas increasing the number of mission-level PDDL actions (\Sfour) leads to a much larger increase, with A* search requiring more than an order of magnitude more time than eager greedy search while both strategies produce plans of similar length. This disparity becomes even more pronounced when all dimensions are combined (\Sfive), where A* search becomes impractical for online TACA, while eager greedy search remains manageable. Peak memory usage grows mainly with map size and, for A* search, with the number of mission-level PDDL actions. It ranges from $37.69$~MiB to $104.66$~MiB in \Sone and reaches $308.43$~MiB for A* search and $90.32$~MiB for eager greedy search in the largest \Sfive instances. \Cref{planta:tab:scalability_summary} summarizes the results, and \Cref{planta:fig:scalability_time} shows the planning-time trends of \Sone, \Sfour, and \Sfive. The plots for all five experiments, including plan length, are given in \Crefrange{planta:S-fig:s1_results}{planta:S-fig:s5_results} of the electronic supplement.

\begin{table}[tb]
\centering
\caption{Scalability experiment results at the maximum tested parameter value for each experiment (\Sone: 900 nodes, \Stwo: 15 architectural configurations, \Sthree: 4 corridor types, \Sfour: 5 mission-level actions, and \Sfive: 90 nodes).}\label{planta:tab:scalability_summary}
\resizebox{.75\textwidth}{!}{%
\begin{tabular}{lcccccccc}
\toprule
\multirow{3}{*}{\textbf{Exp.}} & \multicolumn{4}{c}{\textbf{A*}} & \multicolumn{4}{c}{\textbf{Eager greedy}} \\
\cmidrule(lr){2-5}\cmidrule(lr){6-9}
 & \multicolumn{2}{c}{\textbf{Planning time (s)}} & \multicolumn{2}{c}{\textbf{Plan length}} & \multicolumn{2}{c}{\textbf{Planning time (s)}} & \multicolumn{2}{c}{\textbf{Plan length}} \\
\cmidrule(lr){2-3}\cmidrule(lr){4-5}\cmidrule(lr){6-7}\cmidrule(lr){8-9}
 & \textbf{Mean} & \textbf{Std.} & \textbf{Mean} & \textbf{Std.} & \textbf{Mean} & \textbf{Std.} & \textbf{Mean} & \textbf{Std.} \\
\midrule
\Sone & 0.92 & 0.02 & 47.20 & 7.74 & 0.95 & 0.03 & 47.20 & 7.74 \\
\Stwo & 0.30 & 0.00 & 20.70 & 1.42 & 0.30 & 0.02 & 20.70 & 1.42 \\
\Sthree & 0.20 & 0.00 & 20.30 & 2.26 & 0.20 & 0.00 & 20.30 & 2.26 \\
\Sfour & 15.45 & 1.24 & 46.00 & 5.52 & 1.05 & 0.30 & 46.80 & 6.21 \\
\Sfive & 188.96 & 20.43 & 48.10 & 3.73 & 13.73 & 3.66 & 51.60 & 5.64 \\
\bottomrule
\end{tabular}%
}
\end{table}


\begin{figure}[tbp]
    \centering
    % Compact legends for the generated plots, so that the legend of (a) does not cover its curves.
    \pgfplotsset{every axis post/.append style={legend style={inner sep=1pt,row sep=-2pt}}}%
    \begin{subfigure}[t]{0.32\linewidth}
        \centering
        {% Generated plot: thin the x ticks so they do not overlap at this panel width.
        \pgfplotsset{every axis post/.append style={xtick={0,300,600,900}}}%
        % Generated by navigation_planta; do not edit by hand.
\begin{tikzpicture}[baseline=(paperAxis.north)]
\begin{axis}[name=paperAxis,
width=\textwidth,
height=0.75\textwidth,
xlabel={Number of Nodes},
ylabel={Planning Time [seconds]},
xmin=0,
xmax=900,
ymin=0,
ymax=1,
xtick={0,200,400,600,800,900},
grid=major,
legend pos=north west,
legend style={font=\footnotesize,nodes={scale=0.8,transform shape}},
legend image post style={scale=0.5}
]
\addplot[name path=upper0,draw=none,forget plot] table[x=nodes,y=mean_time_plus_std,col sep=comma] {pddl/navigation_experiment/data/s1_a_star_experiment_processed.csv};
\addplot[name path=lower0,draw=none,forget plot] table[x=nodes,y=mean_time_minus_std,col sep=comma] {pddl/navigation_experiment/data/s1_a_star_experiment_processed.csv};
\addplot[fill=\greenColor!20,fill opacity=0.3,draw=none,forget plot] fill between[of=upper0 and lower0];
\addplot[thick,color=\greenColor] table[x=nodes,y=mean_time,col sep=comma] {pddl/navigation_experiment/data/s1_a_star_experiment_processed.csv};
\addlegendentry{A* mean}
\addplot[name path=upper1,draw=none,forget plot] table[x=nodes,y=mean_time_plus_std,col sep=comma] {pddl/navigation_experiment/data/s1_eager_experiment_processed.csv};
\addplot[name path=lower1,draw=none,forget plot] table[x=nodes,y=mean_time_minus_std,col sep=comma] {pddl/navigation_experiment/data/s1_eager_experiment_processed.csv};
\addplot[fill=\blueColor!20,fill opacity=0.3,draw=none,forget plot] fill between[of=upper1 and lower1];
\addplot[thick,color=\blueColor] table[x=nodes,y=mean_time,col sep=comma] {pddl/navigation_experiment/data/s1_eager_experiment_processed.csv};
\addlegendentry{EG mean}
\addplot[name path=upper2,draw=none,forget plot] table[x=nodes,y=mean_time_plus_std,col sep=comma] {pddl/navigation_experiment/data/s1_a_star_no_sas_experiment_processed.csv};
\addplot[name path=lower2,draw=none,forget plot] table[x=nodes,y=mean_time_minus_std,col sep=comma] {pddl/navigation_experiment/data/s1_a_star_no_sas_experiment_processed.csv};
\addplot[fill=\orangeColor!20,fill opacity=0.3,draw=none,forget plot] fill between[of=upper2 and lower2];
\addplot[thick,color=\orangeColor] table[x=nodes,y=mean_time,col sep=comma] {pddl/navigation_experiment/data/s1_a_star_no_sas_experiment_processed.csv};
\addlegendentry{No adap. mean}
\end{axis}
\end{tikzpicture}
}
        \caption{\Sone: map size.}%
        \label{planta:plot:s1_planning_time}
    \end{subfigure}\hfill
    \begin{subfigure}[t]{0.32\linewidth}
        \centering
        % Generated by navigation_planta; do not edit by hand.
\begin{tikzpicture}[baseline=(paperAxis.north)]
\begin{axis}[name=paperAxis,
width=\textwidth,
height=0.75\textwidth,
xlabel={Number of mission actions},
ylabel={Planning Time [seconds]},
xmin=1,
xmax=5,
ymin=0,
ymax=18.3531757133,
xtick={1,2,3,4,5},
grid=major,
legend pos=north west,
legend style={font=\footnotesize,nodes={scale=0.8,transform shape}},
legend image post style={scale=0.5}
]
\addplot[name path=upper0,draw=none,forget plot] table[x=m,y=mean_time_plus_std,col sep=comma] {pddl/navigation_experiment/data/s4_astar_experiment_processed.csv};
\addplot[name path=lower0,draw=none,forget plot] table[x=m,y=mean_time_minus_std,col sep=comma] {pddl/navigation_experiment/data/s4_astar_experiment_processed.csv};
\addplot[fill=\greenColor!20,fill opacity=0.3,draw=none,forget plot] fill between[of=upper0 and lower0];
\addplot[thick,color=\greenColor] table[x=m,y=mean_time,col sep=comma] {pddl/navigation_experiment/data/s4_astar_experiment_processed.csv};
\addlegendentry{A* mean}
\addplot[name path=upper1,draw=none,forget plot] table[x=m,y=mean_time_plus_std,col sep=comma] {pddl/navigation_experiment/data/s4_eager_experiment_processed.csv};
\addplot[name path=lower1,draw=none,forget plot] table[x=m,y=mean_time_minus_std,col sep=comma] {pddl/navigation_experiment/data/s4_eager_experiment_processed.csv};
\addplot[fill=\blueColor!20,fill opacity=0.3,draw=none,forget plot] fill between[of=upper1 and lower1];
\addplot[thick,color=\blueColor] table[x=m,y=mean_time,col sep=comma] {pddl/navigation_experiment/data/s4_eager_experiment_processed.csv};
\addlegendentry{EG mean}
\addplot[name path=upper2,draw=none,forget plot] table[x=m,y=mean_time_plus_std,col sep=comma] {pddl/navigation_experiment/data/s4_astar_no_sas_experiment_processed.csv};
\addplot[name path=lower2,draw=none,forget plot] table[x=m,y=mean_time_minus_std,col sep=comma] {pddl/navigation_experiment/data/s4_astar_no_sas_experiment_processed.csv};
\addplot[fill=\orangeColor!20,fill opacity=0.3,draw=none,forget plot] fill between[of=upper2 and lower2];
\addplot[thick,color=\orangeColor] table[x=m,y=mean_time,col sep=comma] {pddl/navigation_experiment/data/s4_astar_no_sas_experiment_processed.csv};
\addlegendentry{No adap. mean}
\end{axis}
\end{tikzpicture}

        \caption{\Sfour: mission-level actions.}%
        \label{planta:plot:s4_planning_time}
    \end{subfigure}\hfill
    \begin{subfigure}[t]{0.32\linewidth}
        \centering
        % Generated by navigation_planta; do not edit by hand.
\begin{tikzpicture}[baseline=(paperAxis.north)]
\begin{axis}[name=paperAxis,
width=\textwidth,
height=0.75\textwidth,
xlabel={Number of nodes},
ylabel={Planning Time [seconds]},
xmin=9,
xmax=90,
ymin=0,
ymax=220,
xtick={10,20,30,40,50,60,70,80,90},
grid=major,
legend pos=north west,
legend style={font=\footnotesize,nodes={scale=0.8,transform shape}},
legend image post style={scale=0.5}
]
\addplot[name path=upper0,draw=none,forget plot] table[x=nodes,y=mean_time_plus_std,col sep=comma] {pddl/navigation_experiment/data/s5_astar_experiment_processed.csv};
\addplot[name path=lower0,draw=none,forget plot] table[x=nodes,y=mean_time_minus_std,col sep=comma] {pddl/navigation_experiment/data/s5_astar_experiment_processed.csv};
\addplot[fill=\greenColor!20,fill opacity=0.3,draw=none,forget plot] fill between[of=upper0 and lower0];
\addplot[thick,color=\greenColor] table[x=nodes,y=mean_time,col sep=comma] {pddl/navigation_experiment/data/s5_astar_experiment_processed.csv};
\addlegendentry{A* mean}
\addplot[name path=upper1,draw=none,forget plot] table[x=nodes,y=mean_time_plus_std,col sep=comma] {pddl/navigation_experiment/data/s5_eager_experiment_processed.csv};
\addplot[name path=lower1,draw=none,forget plot] table[x=nodes,y=mean_time_minus_std,col sep=comma] {pddl/navigation_experiment/data/s5_eager_experiment_processed.csv};
\addplot[fill=\blueColor!20,fill opacity=0.3,draw=none,forget plot] fill between[of=upper1 and lower1];
\addplot[thick,color=\blueColor] table[x=nodes,y=mean_time,col sep=comma] {pddl/navigation_experiment/data/s5_eager_experiment_processed.csv};
\addlegendentry{EG mean}
\end{axis}
\end{tikzpicture}

        \caption{\Sfive: combined stress test.}%
        \label{planta:plot:s5_planning_time}
    \end{subfigure}
    \caption{Planning time in the scalability experiments \Sone, \Sfour, and \Sfive. The shaded area shows the standard deviation. EG stands for the eager greedy search.}%
    \label{planta:fig:scalability_time}
\Description{Three line-plot panels show mean planning time in seconds. In S1, as map size increases
from 9 to 900 retained nodes, A* and eager greedy search have similar planning times, reaching
0.92 and 0.95 seconds, respectively, and the no-adaptation baseline reaches about 0.15 seconds. In
S4, as the number of mission-level action schemas increases from 1 to 5, A* planning time rises
sharply to 15.45 seconds, while eager greedy search reaches 1.05 seconds and the no-adaptation
baseline remains near 0.10 seconds. In S5, the combined stress test, as map size increases from 9
to 90 retained nodes, A* planning time rises to 188.96 seconds, whereas eager greedy search reaches
13.73 seconds. Shaded bands indicate standard deviations.}
\end{figure}

\emph{\Sone: Map size.}
\Cref{planta:plot:s1_planning_time} shows that the planning time increases with map size. For the largest maps, with $900$ retained nodes, computing a plan takes less than one second with both search strategies (\Cref{planta:tab:scalability_summary}). This is adequate for non-safety-critical robotic applications in which replanning is not required at high frequency, such as the experiment's navigation application. The A* and eager greedy curves are similar, indicating little difference between optimal and satisficing planning in this experiment. A* and eager greedy also result in plans with the same length (\Cref{planta:tab:scalability_summary}).

\emph{\Stwo and \Sthree: Architectural configurations and quality requirements.}
Planning time remains below $0.31$ seconds over the evaluated range of architectural configurations and corridor types (\Cref{planta:tab:scalability_summary}). A* and eager greedy return plans with the same length.

\emph{\Sfour: Mission-level PDDL actions.}
\Cref{planta:plot:s4_planning_time} shows that planning time increases with the number of mission-level PDDL actions, especially for A* search. With $5$ mission-level actions, A* requires $15.45$ seconds on average, whereas eager greedy search requires $1.05$ seconds. \Cref{planta:tab:scalability_summary} also shows that this reduction in planning time does not come with a substantial increase in plan length. This latency matters for online TACA\@: in the navigation scenario, for example, a lamp failure before traversing a dark corridor requires replanning so that the resulting task plan and architectural configuration satisfy the corridor requirements. Therefore, eager greedy search is the more practical choice for online TACA in this scenario.

\emph{\Sfive: Combined stress test.}
When the dimensions varied individually across \Sone--\Sfour are combined, planning time increases substantially compared to varying map size alone. In \Sfive, planning time reaches $188.96$ seconds on average for A* search and $13.73$ seconds on average for eager greedy search at just $90$ retained nodes, as shown in \Cref{planta:plot:s5_planning_time}. This represents an increase of $205.62$ times in A* planning time relative to \Sone, for a map that is one-tenth the size~(90 vs 900 retained nodes), confirming that combining all dimensions simultaneously amplifies planning time well beyond what any single-dimension experiment produces. Despite the difference in planning time between both search strategies, plan length has a modest difference (\Cref{planta:tab:scalability_summary}). For online TACA under the full-complexity configuration, eager greedy search is the more tractable option. Nevertheless, its planning time is still high for applications requiring frequent replanning.

\emph{No adaptation.}
\Cref{planta:plot:s1_planning_time} and \Cref{planta:plot:s4_planning_time}\footnote{The no-adaptation curve in both plots reports A* search. The corresponding eager greedy curve is omitted because it follows the same trend and does not add relevant information.} demonstrate that planning with only the mission logic~(i.e., task planning) is substantially faster than planning with the adaptation logic included~(i.e., task and architecture planning). For the single-action mission used in \Sone--\Sthree, the no-adaptation baseline reaches approximately $0.15$ seconds at $900$ retained nodes. For the mission used in \Sfour--\Sfive, the no-adaptation baseline reaches approximately $0.10$ seconds with five mission-level actions. These results indicate that the adaptation logic is the primary source of the planning-time increase observed with adaptation.

\emph{Answer to RQ3:} The scalability experiments \Sone--\Sthree show that planning time remains manageable as map size, the number of architectural configurations, and the number of quality requirements increase individually, for both satisficing and optimal planning. \Sfour shows that increasing the number of mission-level PDDL actions has a substantially greater impact on planning time under optimal planning than under satisficing planning. Since both approaches produce plans of similar length, satisficing planning is the more suitable choice for this scenario. \Sfive shows that, when all dimensions are combined, planning time reaches impractical levels under optimal planning and high, though potentially acceptable, levels under satisficing planning for non-safety-critical applications with low replanning frequency. Across all experiments, \Cref{planta:tab:pddl_instances} shows that raw instance size, in objects and initial facts, is driven primarily by map size rather than by architectural or mission complexity, indicating that the planning-time increase observed in \Sfour and \Sfive stems from the added branching factor of additional PDDL action schemas rather than from a larger planning instance.
% Based on these results, a practical strategy is to use optimal planning for the initial mission plan and satisficing planning for runtime replanning.

%%% Local Variables:
%%% mode: latex
%%% TeX-master: "main"
%%% End:
